\clearpage\begingroup\NoSourceHead\SourceLeaf{206}
\begin{center}
{[II.]}\\[.8em]
\textbf{BERICHTIGUNG ZU DEM AUFSATZE:}\\[.4em]
\textbf{BERECHNUNG DES OSTERFESTES MON. CORR. 1800 AUG. S. 121[*)].}\\[.7em]
\rule{\textwidth}{.35pt}\\[.2em]
{\small Zeitschrift für Astronomie und verwandte Wissenschaften, herausgeg. von v. \textsc{Lindenau} u. \textsc{Bohnenberger},\\
1, 1816 Jan. u. Febr., S. 158.}\\[.2em]
\rule{\textwidth}{.35pt}
\end{center}
\begingroup\renewcommand{\thefootnote}{}\footnotetext{[*) Werke VI, S. 73.]}\endgroup
\noindent\hspace*{\parindent}Die Hülfsgrösse $M$ wird durch die S. 129 [Werke VI, S. 78] gegebene Regel vom Jahre 4200 bis 4300, von 4500 bis 4600, von 4800 bis 4900 u. s. f. um eine Einheit zu klein gefunden, und in spätern Jahrhundert[en] kann der Unterschied zwischen dem Resultate jener Regel und dem, was die Einrichtung des gregorianischen Kalenders erfordert, immer grösser werden. Ich verdanke diese Bemerkung dem Hrn. Dr. \textsc{Tittel} aus Erlau, welcher sich gegenwärtig in Göttingen aufhält, und mit ausgezeichnetem Eifer dem Studium der astronomischen Wissenschaften widmet.

Der Grund dieses Unterschieds liegt in dem von mir bey Aufstellung jener Regel nicht beachteten Umstande, dass die Urheber des gregorianischen Kalenders die sogenannte Mondsgleichung, die bis 3900 incl. aller 300 Jahre angebracht wird, von 4200 auf 4300 verlegen, weil ihrer Meinung nach diese Gleichung eigentlich aller $312\frac12$ Jahr einmal angebracht werden muss. Ohne mich hier auf die Frage einzulassen, ob durch diese Anordnung der beabsichtigte Zweck erreicht wird, bemerke ich nur, dass meine Regel mit der Einrichtung des gregorianischen Kalenders, wie er ist, leicht in völlige Übereinstimmung gebracht werden kann, wenn nur für die Zahl $p$ nicht wie a. a. O. vorgeschrieben ist, der Quotient bey Division der Zahl $k$ mit 3, sondern der Quotient bey Division der Zahl $13+8k$ mit 25 angenommen wird.
\HistoricalFoot{XI\textsubscript{1}.}{26}
\clearpage\SourceLeaf{207}\HistoricalHead{202}{L}{CHRONOLOGIE.\quad VERÖFFENTLICHUNGEN.}
\begin{center}\textbf{BEMERKUNGEN.}\end{center}
\begingroup\small
\noindent\hspace*{\parindent}Daß \textsc{Gauss}' ursprüngliche Vorschrift (Monatliche Correspondenz der Erd- und Himmels-Kunde 2, August 1800, S. 128, Werke VI, S. 78) über die Bestimmung der Zahl $p$ vom Jahre 4200 an ihre Gültigkeit verliert, hat vor \textsc{Gauss} als erster \textsc{Français}*) öffentlich festgestellt und gleichzeitig die Zahl $p$ durch die Formel $p=Q\left(\dfrac{k-Q\left(\dfrac{k-17}{25}\right)}{3}\right)$ bestimmt, die unbeschränkte Gültigkeit hat; hierbei soll allgemein $Q\left(\dfrac{a}{b}\right)$ den ganzzahligen Quotienten bedeuten, der sich bei Division der ganzen Zahl $a$ durch die ganze Zahl $b$ ergibt.\begingroup\renewcommand{\thefootnote}{}\footnotetext{*) J. F. \textsc{Français}, \textit{Solution directe des principaux problèmes du calendrier}, Annales de Mathématiques pures et appliquées, 4, 1813, S. 273.}\endgroup\ Für die Rechnung ist \textsc{Français}' Formel etwas bequemer als die \textsc{Gauss}sche Festsetzung $p=Q\left(\dfrac{8k+13}{25}\right)$, die \textsc{Gauss} auch handschriftlich in sein Exemplar der Monatlichen Correspondenz 2, S. 128 eingetragen hat (Vergl. Werke VI, S. 79, Z. 5 v. u.). Weiter ist vor \textsc{Gauss}' Veröffentlichung \textsc{Delambre} in der oben S. 200, dritte Fußnote genannten Abhandlung (vergl. besonders a. a. O. S. 316), anscheinend ohne \textsc{Français}' Aufsatz zu kennen, auf die gleiche Verbesserung wie dieser gekommen. Vergl. ferner \textsc{Delambre}, \textit{Histoire de l'Astronomie moderne}, t. I, Paris 1821, S. 25. Auch L. \textsc{Ciccolini}**) hat nach seinen Angaben seit 1813 erkannt, daß die \textsc{Gauss}sche Osterformel vom Jahre 4200 an eine Änderung erfordert; er hat daher die von $p$ abhängige \textsc{Gauss}sche Zahl $M$ (Werke VI, S. 78) in einer von \textsc{Gauss} abweichenden Weise bestimmt.\begingroup\renewcommand{\thefootnote}{}\footnotetext{**) L. \textsc{Ciccolini}, \textit{Formole analitiche pel calcolo della Pasqua e correzione di quelle di Gauss con critiche osservazioni su quanto ha scritto del calendario il Delambre}, Roma 1817, vergl. Vorrede S. VII, sowie S. 76. Nachtrag hierzu: \textit{Al signor Francesco Carlini astronomo di Brera}, Biblioteca italiana 13, 1819, S. 350; vergl. auch die S. 346 voraufgehende Besprechung des \textsc{Ciccolini}schen Buches, als deren Verfasser \textsc{Carlini} anzusehen ist.}\endgroup\ Schließlich hat P. \textsc{Tittel}, der die vorstehende \textit{Berichtigung} veranlaßt hat, in seiner Arbeit \textit{Methodus technica, brevis, perfacilis ac perpetua construendi calendarium ecclesiasticum stylo tam novo quam vetere}, Goettingae 1816, § 11, Solutio I, S. 10***) bei der sogenannten Mondgleichung den Wert $Q\left(\dfrac{8k+13}{25}\right)$.\begingroup\renewcommand{\thefootnote}{}\footnotetext{***) Diese Stelle gehört nicht zu denjenigen, die \textsc{Tittel} in der Einleitung seiner Schrift als von \textsc{Gauss} stammend bezeichnet; vergl. den folgenden Abschnitt [III.].}\endgroup

Infolge der für die Zahl $p$ notwendig werdenden Korrektur sind auch \textsc{Gauss}' Angaben in seinem Aufsatz \textit{Noch etwas über die Bestimmung des Osterfestes}, Braunschweigisches Magazin 1807 September 12 (Werke VI, S. 86), über diejenigen Jahrhunderte, für die Ostern im julianischen und im gregorianischen Kalender auf das gleiche Datum fällt, unzutreffend. Für den richtigen Wert $p=Q\left(\dfrac{8k+13}{25}\right)$ ergibt sich, daß bereits im 68. Jahrhundert, d. i. vom Jahre 6700 bis 6799, also weit früher als \textsc{Gauss} findet, Ostern in beiden Kalendern auf einerlei Datum fallen würde, wenn dieselben solange ungeändert im Gebrauch bleiben sollten. Einerlei Datum würde aber alsdann in den beiden Kalendern immer um 7 volle Wochen auseinander sein. Mit dem Jahre 6800 hört diese Übereinstimmung wieder auf und tritt von 20700 bis 20799 von neuem ein, wo aber der Unterschied 22 Wochen betragen würde. In diesem Sinne (vgl. \textsc{Franz Goldscheider}, \textit{Über die Gausssche Osterformel}, II. Teil, Wissenschaftliche Beilage zum Jahresbericht des Luisenstädtischen Realgymnasiums zu Berlin, Ostern 1899, S. 6 und 7) ist der Inhalt des vorletzten Abschnitts der angeführten \textsc{Gauss}schen Arbeit zu berichtigen.
\begin{flushright}\textsc{Loewy.}\end{flushright}
\endgroup
\begin{center}\rule{.24\textwidth}{.35pt}\end{center}
\clearpage\endgroup
